% H200 tuned ADR campaign versus the RTX PRO 5000 Blackwell results of this study.
% Tables and figures: scripts/reports/make_adr_h200_comparison.py (no solver execution);
% inputs: campaigns/adr_h200_tuned/records (export of 2026-10-29) and this study's CSV files.
\section{H200 with tuned kernels versus the RTX PRO 5000 Blackwell}
\label{sec:adr-h200-tuned}

The H200 tuned campaign (\texttt{campaigns/adr\_h200\_tuned}, 28--29 October 2026) solved the
systems of this study again on one NVIDIA H200 NVL (141~GiB, FP64 peak estimate 30.2~TFLOP/s from
the campaign preflight), then continued the mesh ladder to about 20~million trace unknowns. The
Blackwell results above were measured on an RTX PRO 5000 Blackwell workstation GPU with about
48~GiB. Three runs enter the comparison:
\begin{itemize}
\item \textbf{Blackwell}: this study (smooth, oscillatory, focused and native-completion suites);
\item \textbf{H200, stock kernels}: the unified campaign 21958 on the same H200, with the
  configurations and kernels of this study;
\item \textbf{H200, tuned kernels}: interleaved face-dense operator, FP32 interleaved ASM
  application, FP64 interleaved block Jacobi, $p$MG--AMG with graph-captured kernels, the
  true-target GMRES restart trigger, and the whole system resident on the device from assembly
  to the last solve.
\end{itemize}

\paragraph{Matching.} A Blackwell measurement and an H200 job are compared only when geometry,
case, triangle count and degree agree and the H200 candidate's recorded alias names the
Blackwell candidate; the H200 plan records the archived operator hash of each of the study's
systems as its identity. The H200 reassembles those systems with its own kernels, so its
matrices are the same discrete systems but hash differently (all 27 checked); the stock and tuned
H200 assemblies, compared directly, agree to $6\cdot10^{-12}$ relative.
Of 75 matched configurations, 71 passed on both machines. Three that passed on Blackwell are
recorded as numerical failures on the H200 (BJ+PP on two variable-velocity and one trigonometric
system at $p=6$, see below), and one passed only on the H200. Times are the medians of the
measured samples; ``repeated'' is the second solve on the same setup and ``amortized'' the setup
plus both solves, per solve, as in the study's CSV files.

\subsection{Same configuration, same operator}

Table~\ref{tab:adr-h200-summary} gives the Blackwell-to-H200 time ratio per family. With the
study's own kernels the H200 is only 1.1--1.5 times faster: these solvers are bound by memory
traffic and kernel launches at the study's sizes, and the H200's advantage in bandwidth and FP64
rate is not used by kernels tuned for the workstation GPU. With the tuned kernels the ratio
grows with the system size (Figure~\ref{fig:adr-h200-speedup}): from 1.0--3.4 below $10^5$
trace DOFs to 4.0--7.5 for ASM+PP, 2.9--7.9 for BJ+PP and 2.2--4.1 for $p$MG--AMG at
$0.34$--$1.04\cdot10^6$ DOFs. AMGX stays at 1.1--1.3.

\begin{table}[ht]
\centering\small
\resizebox{\linewidth}{!}{\input{\ADRResultsPath/h200_tuned/summary_table.tex}}
\caption{Blackwell time divided by H200 time on matched configurations: median and range.
Stock: unified campaign 21958 on the H200 (27 matched configurations); tuned: this campaign.
Per-iteration: repeated time per iteration, which removes the effect of a different iteration
count.}
\label{tab:adr-h200-summary}
\end{table}

\begin{figure}[ht]
\centering
\includegraphics[width=\linewidth]{\ADRResultsFigurePath/h200_tuned/speedup_matched.pdf}
\caption{Blackwell time over H200 time for each matched configuration, against trace DOFs.}
\label{fig:adr-h200-speedup}
\end{figure}

Two effects make up the tuned speed-up, and the table separates them:
\begin{itemize}
\item \textbf{Cost per iteration.} At about $10^6$ DOFs one iteration is 4.6--5.1 times cheaper
  for ASM+PP, 4.0 for BJ+PP, 2.7--2.8 for $p$MG--AMG and 1.3 for AMGX. Below $10^5$ DOFs the
  factor is 1.0--1.9: kernel launch latency dominates there on both GPUs.
\item \textbf{Iteration count.} The true-target trigger stops GMRES when the true residual meets
  the target instead of over-solving to the estimated one. On 39 of the 71 configurations the
  H200 needs fewer iterations, on 29 the same and on 3 slightly more. ASM+PP and $p$MG--AMG
  often need half as many on the trigonometric case (for example 126 versus 64 at
  $3.4\cdot10^5$ DOFs). The three BJ+PP rows recorded as numerical failures on the H200 met the
  residual target but missed the study's $L^2$ rule.
\end{itemize}

AMGX gains least. Its setup is 1.6--7.5 times faster on the H200 (median 3.6): the matrix is
built on the device and handed over device to device. Its V-cycles are the same library code on
both GPUs and improve by 1.1--1.3 per iteration. Setup falls most for the block preconditioners,
from 12--1{,}390~ms on Blackwell to 4--56~ms at the matched sizes, because the local
factorizations and the operator layout are now built on the device.

\subsection{Did the winner change?}

Table~\ref{tab:adr-h200-family} compares the fastest passing member of each family on each
machine, for the 25 systems both measured, against the confirmed PyPardiso LU baseline of the
study (measured on the Blackwell workstation's host CPU; the LU campaign on a genoa CPU node of
the H200's cluster is reported in Section~\ref{sec:adr-h200-node-lu}).
\begin{itemize}
\item \textbf{AMGX no longer wins.} On Blackwell AMGX had the fastest repeated solve on the
  largest smooth systems (transport, anisotropic and trigonometric at 99{,}458 triangles,
  $p=6$). On the H200 $p$MG--AMG wins the anisotropic, trigonometric and variable-velocity
  systems from 32{,}768 triangles up, and the block preconditioners win transport; AMGX wins
  none of the 25.
\item \textbf{The crossover against LU moves to smaller meshes.} For repeated solves at $p=6$
  the best iterative solver beats the reused LU factors on transport and variable velocity from
  8{,}192 triangles on both machines. On the trigonometric case the H200 crosses at 8{,}192
  triangles and Blackwell at 32{,}768; on the anisotropic case at 32{,}768 and 99{,}458. At
  99{,}458 triangles the best iterative time over the LU time is 0.07--0.21 on the H200 against
  0.19--0.82 on Blackwell for the first three cases, and 0.55 against 0.90 for the anisotropic
  one.
\item \textbf{Cellular annulus: LU still wins repeated solves, the H200 wins fresh ones.} On
  \texttt{cellular7\_weak} the reused LU factors stay 2--16 times faster than any iterative
  solve on either machine. For a fresh solve (setup included) the H200's ASM+PP beats LU from
  8{,}039 triangles up (0.41--0.52 of the LU time), where Blackwell needed 1.4--3.2 times the
  LU time.
\item \textbf{Below about 8{,}000 triangles LU keeps the repeated solve}: a reused LU solve
  takes about 5~ms and every iterative method on either GPU is slower. A fresh solve is another
  matter: there the H200's best iterative solver needs 0.23--0.49 of the LU time on the smooth
  cases already at 2{,}048 triangles.
\end{itemize}

\begin{table}[ht]
\centering\scriptsize
\input{\ADRResultsPath/h200_tuned/family_table.tex}
\caption{Fastest passing family per system for the repeated solve (ms), and the study's
P-LU reused solve. ``BW'': Blackwell.}
\label{tab:adr-h200-family}
\end{table}

\subsection{H200 versus PyPardiso LU on a genoa CPU node}
\label{sec:adr-h200-node-lu}

The PARDISO v2 campaign factored the exported systems of the unified device campaign with
PyPardiso (oneMKL) on a CPU node of the H200's cluster (genoa partition, two AMD EPYC 9655, 192
cores, 754~GiB; jobs 24478 and 24482), not on the H200 node's own host CPUs (two 64-core
sockets), at the same residual target ($10^{-10}$). For each system it kept the fastest of the
measured thread counts (8--192), separately for the fresh solve (CSR conversion, analysis, factorization and first
solve) and for a solve that reuses the factors. A follow-up run of the same worker (jobs 26536,
26551 and 26578, records \texttt{h200\_tuned\_lu}) factorized the 68 H200 systems the first plan
lacked, up to L7. All 114 H200 systems therefore have an LU measurement.
Tables~\ref{tab:adr-h200-node-lu} and \ref{tab:adr-h200-node-lu-large} put these next to the
H200's fastest passing iterative solver, fresh (setup and first solve) and repeated. 104 systems
have a passed iterative solve. Ten at L6--L7 have none (seven stopped at the FP64 residual
floor, FP64$^*$; three not passed), and they are not counted as LU wins. Table~\ref{tab:adr-h200-node-lu-crossover}
gives the crossover per case and degree.
\begin{itemize}
\item \textbf{Fresh solves: the GPU wins on 89 of the 104.} On the smooth cases the best iterative
  fresh time is 0.02--0.18 of the LU time at L5--L7 (variable velocity 0.02--0.06). LU keeps the
  fresh solve in two situations:
  \begin{itemize}
  \item the smallest meshes: L1 at $p=2$ (anisotropic, cellular, trigonometric, variable
    velocity), cellular at L2, $p=2$, and at L1, $p=6$;
  \item the hard cases at scale: the cellular annulus at L5--L7 (1.5--2.2 times faster than
    the GPU) and closed-loop stress at L6, $p=2$, and L7 (1.2--3.1 times).
  \end{itemize}
\item \textbf{Repeated solves: the GPU wins on 42 of the 104, all but two of them smooth.}
  Reusing the factors costs 3~ms at about 36{,}000 trace DOFs, 0.23--0.25~s at a million and
  5.0--10.6~s at 20 million. The iterative solver wins:
  \begin{itemize}
  \item transport on 15 of 20 systems, variable velocity on 14 of 20, trigonometric on 11 of 18;
  \item at $p=6$ from 85{,}120 (transport) and 342{,}272 DOFs (trigonometric, variable velocity),
    at $p=9$ from 30{,}080 and 121{,}600, and at $p=2$ only from 446{,}223 (variable velocity)
    or 8~million DOFs (transport, trigonometric);
  \item the anisotropic case only at L5, $p=6$ and 9 (ratio 0.54 and 0.96).
  \end{itemize}
  At the largest size the ratio is 0.06--0.10 for variable velocity ($p$MG--AMG), 0.29--0.36 for
  trigonometric and 0.33--0.71 for transport (AMGX). On the cellular annulus and on closed-loop
  stress the reused factors stay faster at every size up to L7: 3.5--35 and 1.3--26 times.
\item \textbf{Consistent with the Blackwell comparison, with a stronger LU.} On the 23 systems
  that both LU campaigns measured, the genoa node's PyPardiso is 1.1--1.9 times faster than the
  Blackwell workstation's host for a fresh solve and 1.5--3.2 times faster for a reused one
  (most at the larger sizes). The direction of every statement in the previous subsection holds,
  but the repeated-solve crossovers move to larger meshes: trigonometric and variable velocity at
  $p=6$ now cross at 32{,}768 triangles instead of 8{,}192, the anisotropic case no longer
  crosses up to 99{,}458 triangles (ratio 1.70), and at 99{,}458 triangles the ratio is
  0.22--0.27 (trigonometric, variable velocity) rather than 0.07--0.21; transport still crosses
  at 8{,}192. For fresh solves the cellular ratio at $p=6$ is 0.60--0.64 instead of 0.41--0.52,
  and at 2{,}048 triangles and $p=2$ LU now keeps the fresh solve on the trigonometric and
  variable-velocity cases.
\item \textbf{LU up to 20 million DOFs.} Fresh LU takes 2.7--20~s at L5, 14--27~s at L6 and
  36--53~s at L7, with a peak resident memory of at most 52, 70 and 141~GiB (a genoa node has
  754~GiB); 96 and 192 threads are within a few per cent. PyPardiso's 32-bit MKL interface
  limits the matrix's own indices (fewer than $2^{31}$ entries; the L7 matrices have at most
  about $7\cdot10^8$), not the factor. Factors above $2^{31}$ entries factorize and solve; only
  MKL's reported factor count overflows. Three L7 LU solves (anisotropic $p=6$, trigonometric
  $p=2$ and 6) end at a true relative residual of $1.10$--$1.36\cdot10^{-10}$, just above the
  target. That is the same FP64 floor the iterative solvers meet there. Their timings are kept
  and marked $\dagger$.
\end{itemize}

\begin{table}[ht]
\centering\scriptsize
\resizebox{\linewidth}{!}{\input{\ADRResultsPath/h200_tuned/node_lu_table.tex}}
\caption{Fastest passing iterative solver on the H200 against PyPardiso LU on a genoa CPU
node, L1--L4 (ms; $t$: threads of the fastest LU measurement). The ratios are iterative over
LU; below one the iterative solver is faster.}
\label{tab:adr-h200-node-lu}
\end{table}

\begin{table}[ht]
\centering\scriptsize
\resizebox{\linewidth}{!}{\input{\ADRResultsPath/h200_tuned/node_lu_large_table.tex}}
\caption{As Table~\ref{tab:adr-h200-node-lu}, L4b--L7. $\dagger$: LU solve at the residual
floor ($1.10$--$1.36\cdot10^{-10}$). FP64$^*$: no iterative family passed, but a run stopped
at the FP64 residual floor with the best solution's accuracy (Table~\ref{tab:adr-h200-outcomes}).}
\label{tab:adr-h200-node-lu-large}
\end{table}

\begin{table}[ht]
\centering\scriptsize
\resizebox{\linewidth}{!}{\input{\ADRResultsPath/h200_tuned/node_lu_crossover_table.tex}}
\caption{Crossover against the genoa LU per case and degree: the smallest measured trace
DOFs from which the best iterative solver is faster at that and every larger measured size
(``--'': not reached), and the ratio at the largest measured size.}
\label{tab:adr-h200-node-lu-crossover}
\end{table}

\subsection{Beyond the Blackwell's reach}

The H200 ladder continues to L5 (0.4--1.4~million triangles; 1.8--6.0~million trace DOFs),
L6 (8~million) and L7 (20~million). Figure~\ref{fig:adr-h200-ladder} follows one class, the
trigonometric case, through the ladder; Table~\ref{tab:adr-h200-extension} gives the medians
per level, degree and family over the passed jobs of all classes. At these sizes AMGX is the
fastest family wherever it passes: its iteration count stays at 15--27 from L2 to L7, while
$p$MG--AMG needs 390--520 iterations from L5 on and the block preconditioners keep growing, up
to about 800 for ASM+PP and 1{,}600 for BJ+PP at L7. L7 at $p=9$ and all of L8--L9 exceed one H200's memory (device estimates of
138--275~GiB against a 125~GiB budget; 780~GiB of host memory at L8, $p=9$) and were not run.
Over two solves per matrix (the amortized measure), the genoa LU is faster than the best H200
solver only on the smallest $p=2$ mesh (8.6 against 17.9~ms). On every other trigonometric
system where both were measured it is 1.2 to 12.7 times slower: at L4, $p=6$, $p$MG--AMG needs
91~ms against the LU's 1{,}152~ms.

\begin{figure}[ht]
\centering
\includegraphics[width=\linewidth]{\ADRResultsFigurePath/h200_tuned/h200_ladder_amortized.pdf}
\caption{Amortized solve time (setup and two solves, per solve) along the ladder for the
trigonometric case: H200 iterative families and PyPardiso LU on a genoa node (dashed; fresh and
reused each at its fastest thread count). Missing iterative points are jobs that did not pass;
hollow markers on a dotted line are those that stopped at the FP64 residual floor with the best
accuracy, at the time they took to reach it (NC$^*$, paragraph on the FP64 floor below). The LU
point at L7 is at the residual floor too. The last points at $p=2$ (34 million DOFs) and $p=9$ (14
million) are the Lmax systems, which have no LU yet; L7 does not fit one H200 at $p=9$.}
\label{fig:adr-h200-ladder}
\end{figure}

\paragraph{Outcomes at scale: the FP64 residual floor ($^*$).}
\label{par:adr-h200-fp64}
A job passes when the solver reports convergence to its internal target ($10^{-11}$) and the
independently recomputed true relative residual is at most $10^{-10}$. A computed residual
$b-Ax$ carries rounding of about $\varepsilon\,\||A||x|\|$, so the relative residual cannot be
evaluated below $\varepsilon\,\||A||x|\|_2/\|b\|_2$, $\varepsilon=2.2\cdot10^{-16}$. This floor was
measured on 39 small exported systems
(\nolinkurl{scripts/advection_diffusion_reaction/diagnostics/fp64_residual_floor.py},
\nolinkurl{h200_tuned/fp64_floor.csv}) and fitted per class as $c\,N^\alpha$ in the trace DOFs $N$:
\begin{itemize}
\item \textbf{Trigonometric and anisotropic:} $\alpha\approx1$, $1.2$--$1.6\cdot10^{-11}$ at one
  million DOFs, $2$--$3\cdot10^{-10}$ at 20 million and about $8\cdot10^{-10}$ at the 34 million
  of Lmax, $p=2$. It passes the internal target at L4--L5 and the acceptance itself at L7, where
  even LU ends at $1.1$--$1.4\cdot10^{-10}$. Their exact
  solution $\sin\pi x\sin\pi y$ vanishes on $\partial\Omega$, so $b$ carries only the source,
  of size $O(h^2)$ per DOF, while $K\approx1$ keeps $|A|$ of order one: the ratio grows like
  $h^{-2}\propto N$.
\item \textbf{The other classes:} $10^{-14}$ to $4\cdot10^{-13}$ at 20 million DOFs. Variable
  velocity has the same $K=1$ but $u_\star\ge1$ on the boundary ($\|b\|\approx10^2$); transport
  has $K=10^{-3}$.
\end{itemize}
An asterisk (NC$^*$, FP64$^*$) marks a job that did not pass, whose true residual is within ten
times its system's floor, and whose $L^2$ error is within 25\,\% of the best on that system or no
worse than that of a solve that passed there. At $p\ge6$ that error is itself round-off
($10^{-12}$ to $10^{-11}$) and varies by tens of percent between runs. 83 jobs qualify, all
trigonometric or anisotropic: 44 at L4--L7 (4 of them from the L4 degree fill) and 39 at Lmax
(Table~\ref{tab:adr-h200-outcomes}):
\begin{itemize}
\item 34 $p$MG--AMG runs iterate on, mostly to the 2{,}000-iteration cap, chasing the internal
  target;
\item 27 ASM+PP runs stop when GMRES detects stagnation or at the cap, many with a true residual
  already below $10^{-10}$;
\item 21 AMGX runs at L6, L7 and Lmax report success to AMGX itself (16--89 iterations), but the
  recomputed residual is $1.2$--$6.6\cdot10^{-10}$;
\item 1 BJ+PP run.
\end{itemize}
\textbf{Time to the floor.} The worker stops a job at its first failed solve, so a starred job
has no timed sample, only its first (warm-up) solve. Its residual history (the true residual at
every restart-cycle end) shows where it first came within a factor of two of its lowest residual;
the graphs plot the setup plus that fraction of the solve time (iterations to that cycle over all
iterations), as a hollow marker, and the same solve fraction as its reused time. The iteration
count is exact for 65 of the 83 jobs and an upper bound for 18, whose restart cycles cannot all be
separated; AMGX stops by itself, so its whole solve counts. Above one million DOFs, warm-up times
match the measured ones: fresh to a median factor of 1.02 (90\,\% within 1.26) and the first solve
the repeated one within 1.07 (113 passed ladder jobs). $p$MG--AMG reaches the floor within
46--276 iterations and then iterates on, mostly to the cap: its time to the floor is 6--50\,\% of
the job.
The grids and tables keep NC$^*$ without a time.

The 93 other failures (45 on the ladder, 11 in the L4 degree fill, 37 at Lmax) are genuine: BJ+PP at L6--L7 and Lmax
(mostly $10^{-10}$ to $4\cdot10^{-3}$ after 2{,}000 iterations), AMGX diverging on the cellular
annulus at Lmax, $p=2$, cellular and stress runs that end far above their own floor, and the polynomial families on
variable velocity (BJ+PP from L3, ASM+PP from L4--L5). Those converge to true residuals of
$6$--$10\cdot10^{-12}$, but their $L^2$ error misses the case's $2\cdot10^{-8}$ bound ($10^{-7}$ to
$10^{-6}$ at L5--L7), while AMGX and $p$MG--AMG reach $10^{-9}$ and $10^{-10}$ on the same systems. The boundary data make $\|b\|$ large
($\approx10^2$), so a small relative residual leaves smooth error that a one-level method without a
coarse space does not remove. A normwise backward-error acceptance,
$\|b-Ax\|/(\||A||x|\|+\|b\|)$, would remove the size dependence of the floor.

\begin{table}[ht]
\centering\small
\input{\ADRResultsPath/h200_tuned/outcome_table.tex}
\caption{H200 job outcomes per level and family. ``FP64 floor$^*$'': not passed, but stopped at
the FP64 residual floor with the best solution's accuracy (paragraph above).}
\label{tab:adr-h200-outcomes}
\end{table}

\begin{table}[ht]
\centering\scriptsize
\input{\ADRResultsPath/h200_tuned/extension_table.tex}
\caption{H200, levels beyond the Blackwell study: medians over the passed jobs of all classes
(ms). Different families pass on different classes, so the medians are not a like-for-like
ranking; Table~\ref{tab:adr-h200-outcomes} gives the coverage.}
\label{tab:adr-h200-extension}
\end{table}

\subsection{Assembly}

The device-route assembly that produced the solvers' matrices takes 1.5~s at L6 and 3.6~s at
L7 for $p=6$; the raw-CUDA tensor kernel (a different discretization, kept as a timing
reference) takes 0.31 and 0.77~s. Both stay well below the solve times from L5 on. The raw kernel
also ran at L8 for $p=2$ and $p=6$ (0.39 and 1.5~s), where no solver fits.
